% 三种选样准则共享完全相同的教学点集、坐标范围和孤立点。
\begin{tikzpicture}[x={\dimexpr\linewidth/174\relax},y=1mm,font=\figurefont\footnotesize,text=BookInk]
  \tikzset{
    candidate/.style={circle,draw=BookMuted,fill=white,line width=.7pt,minimum size=3.2mm,inner sep=0pt},
    chosen/.style={circle,draw=FigureModel,fill=FigureModel,line width=1pt,minimum size=3.8mm,inner sep=0pt},
    panel title/.style={font=\figurefont\bfseries\footnotesize,anchor=south},
    criterion/.style={font=\figurefont\scriptsize,text=BookMuted,align=center}
  }
  \def\candidatepoints{5/9,8/13,10/7,12/17,15/11,20/27,24/10,24/31,27/15,29/29,30/8,33/19,36/13,45/34}
  \foreach \shift/\panel/\title in {0/a/预测不确定性,61/b/代表性,122/c/批内多样性}{
    \begin{scope}[xshift=\shift mm]
      \node[panel title] at (26,46) {(\panel)\quad\title};
      \draw[BookRule,line width=.7pt] (2,5) rectangle (50,40);
      \foreach \x/\y in \candidatepoints{\node[candidate] at (\x,\y) {};}
      \node[criterion] at (26,1.5) {\strut 候选为空心，选中样本为实心};
    \end{scope}
  }
  \begin{scope}
    \draw[FigureLoss,dashed,line width=1pt] (22,5)--(22,40);
    \node[chosen] at (20,27) {};
    \node[chosen] at (24,10) {};
    \node[anchor=south west,text=FigureLoss] at (22.5,35.5) {给定边界};
    \node[anchor=north east,text=BookMuted] at (47,33) {孤立点};
  \end{scope}
  \begin{scope}[xshift=61mm]
    \draw[FigureInput!35,line width=.7pt] (10,11) ellipse (8 and 6);
    \draw[FigureInput!35,line width=.7pt] (31,13) ellipse (8 and 7);
    \draw[FigureInput!35,line width=.7pt] (25,30) circle (7);
    \node[chosen] at (10,11) {};
    \node[chosen] at (31,13) {};
    \node[chosen] at (24,31) {};
    \node[anchor=north east,text=BookMuted] at (47,33) {孤立点};
  \end{scope}
  \begin{scope}[xshift=122mm]
    \node[chosen] at (5,9) {};
    \node[chosen] at (29,29) {};
    \node[chosen] at (36,13) {};
    \draw[BookTeal,dotted,line width=.8pt] (5,9)--(29,29)--(36,13)--cycle;
    \draw[FigureLoss,line width=.8pt] (43,32)--(47,36) (43,36)--(47,32);
    \node[anchor=south east,text=BookMuted,align=right] at (48,34) {孤立点\\未必有用};
  \end{scope}
\end{tikzpicture}
